% 原书分章PDF第19—25页：第26.4.2—26.4.3节。
\subsection{情形II：无限支撑集}

现在假设 $W$ 在每个所考虑的 $P$ 下均具有无限支撑集。在这种情形下，g计算参数在我们所考虑的非参数模型中通常不存在二阶梯度。尽管如此，仍可以考虑近似二阶典范梯度，从而得到一个定义良好的2-TMLE，下面予以说明。

假设 $W:=(W(1),\ldots,W(d))$，且 $W_j$ 为实值。逐点定义
\[
D_h^{(2)}(\bar Q,\bar g,Q_W)(o_1,o_2):=
H_h^{(2)}(\bar g,Q_W)(w_1,a_1,w_2,a_2)[y_1-\bar Q(w_1)],
\]
其中还定义了
\[
\begin{aligned}
&H_h^{(2)}(\bar g,Q_W)(w_1,a_1,w_2,a_2)\\
&\qquad:=\frac1{h^d}K\left(\frac{w_1-w_2}{h}\right)
\frac{2a_1}{\bar g(w_1)Q_W(w_1)}\left[1-\frac{a_2}{\bar g(w_1)}\right],
\end{aligned}
\]
$K$ 为具有紧支撑集的多元核函数，$h$ 为正带宽。由此可见，$D_h^{(2)}$ 是上一小节所定义二阶部分典范梯度 $D^{(2)}$ 的核近似。对每个 $w$，下文以 $K_h(w)$ 表示 $h^{-d}K(h^{-1}w)$。Díaz 等（2016）的引理1确立了一些条件，保证TMLE中由于采用这一近似二阶典范梯度而产生的偏差项渐近可忽略。

原则上，计算g计算参数不需要估计 $W$ 的密度函数；估计其分布函数就已足够。但是，从近似二阶典范梯度的形式可以看出，二阶靶向过程显式需要这一密度函数。因此，要构造2-TMLE，首先必须估计 $W$ 的密度。至少可以考虑以下两种方法来处理这一问题：
\begin{enumerate}
\item 在二阶靶向过程中，每当需要时都使用 $Q_{W,0}$ 的一个固定平滑估计 $Q_{W,n}^\circ$，得到 $\bar Q_n^*$ 的靶向估计；但最终计算 $\psi_0$ 的靶向估计时，以经验测度 $Q_{W,n}$ 为基础；
\item 在二阶靶向过程中，以及计算 $\psi_0$ 的靶向估计时，均使用 $Q_{W,0}$ 的同一个平滑估计 $Q_{W,n}^\circ$。
\end{enumerate}
第二种方案需要大量额外工作。因为要使平滑估计量 $Q_{W,n}^\circ$ 适合用于构造 $\psi_0$ 的靶向估计，它本身也必须经过靶向处理。因此，需要将它与 $\bar Q_0$ 的估计一起迭代更新。第一种方案绕过了这一问题，因为经验分布 $Q_{W,n}$ 是NPMLE，已经求解了相关得分方程。基于这一理由，我们建议在实践中采用第一种方案，以下也仅讨论这一方案。

考虑这样一种2-TMLE：在TMLE算法中，使用经验分布 $Q_{W,n}$ 作为 $Q_{W,0}$ 的初始估计量；但在量
\[
H_h^{(2)}(\bar g_n,Q_{W,n}^\circ)
\]
中，使用一个固定的平滑估计量 $Q_{W,n}^\circ$，并利用这一量，构造经过 $\bar Q_0$ 当前估计 $\bar Q_{n,m}$ 的、基于logistic回归的最不利子模型。这恰好是第26.4节2-TMLE的一个例子。这一算法的实现尤其简单，因为一次靶向步骤就足以实现解析收敛。为执行这一次更新，在logistic回归模型
\[
\bar Q_{n,0,h}(\epsilon):=\expit\left[\logit(\bar Q_{n,0})
+\epsilon_1H^{(1)}(\bar g_n)+\epsilon_2\bar H_{h,n}^{(2)}(\bar g_n,Q_{W,n}^\circ)\right]
\]
中，求取 $\epsilon=(\epsilon_1,\epsilon_2)$ 的极大似然估计 $\epsilon_n$，其中 $H^{(1)}(\bar g_n)(o)=a/\bar g_n(w)$，且
\[
\begin{aligned}
&\bar H_{h,n}^{(2)}(\bar g_n,Q_{W,n}^\circ)(o)\\
&\quad=\frac1n\sum_{j=1}^n\frac1{h^d}K\left(\frac{w-W_j}{h}\right)
\frac{2a}{\bar g_n(w)Q_{W,n}^\circ(w)}\left[1-\frac{A_j}{\bar g_n(w)}\right]\\
&\quad=\frac{2a}{\bar g_n(w)Q_{W,n}^\circ(w)}
\left[\frac1n\sum_jK_h(w-W_j)
-\frac{\frac1n\sum_jK_h(w-W_j)A_j}{\bar g_n(w)}\right].
\end{aligned}
\]
若采用核估计 $Q_{W,n}^\circ(w):=n^{-1}\sum_{i=1}^nK_h(w-W_i)$，容易验证，上式可简化为
\[
\bar H_{h,n}^{(2)}(\bar g_n,Q_{W,n}^\circ)(o)=\frac{2a}{\bar g_n(w)}
\left[1-\frac{\bar g_{n,h}(w)}{\bar g_n(w)}\right],
\]
其中 $\bar g_{n,h}(w):=\sum_jK_h(w-W_j)A_j/\sum_jK_h(w-W_j)$ 是由带宽 $h$ 索引的 $\bar g_0(w)$ 的Nadaraya--Watson估计量。

给定通过这一单次更新程序得到的 $\bar Q_0$ 的靶向估计 $\bar Q_{n,h_n}^*:=\bar Q_{n,0,h_n}(\epsilon_n)$，所得 $\psi_0$ 的2-TMLE为
% 原书行内最终估计量改作无编号展示。
\[
\psi_n^*:=\Psi(\bar Q_{n,h_n}^*,Q_{W,n})=\frac1n\sum_{i=1}^n\bar Q_{n,h_n}^*(W_i).
\]
下面还将用到三阶余项 $R_3$，与前面一样，定义为
\[
R_3(P,P_0):=P_0\left[\left(1-\frac{Q_{W,0}\bar g_0}{Q_W\bar g}\right)
\left(\frac{\bar g-\bar g_0}{\bar g}\right)(\bar Q-\bar Q_0)\right].
\]
Díaz 等（2016）的定理1给出了使 $\psi_n^*$ 成为 $\psi_0$ 的渐近线性且有效估计量的条件。该定理包括以下条件：
\begin{enumerate}
\item 核函数 $K$ 为 $2d$ 次可微，且 $h_n^{2d}n\to+\infty$；
\item $R_3(P_n^*,P_0)=o_P(n^{-1/2})$，且 $\|\bar Q_n^*-\bar Q_0\|h_n^{m_0+1}=o_P(n^{-1/2})$；
\end{enumerate}
其中 $m_0$ 使得 $\bar g_0$ 和 $\bar Q_0$ 都为 $(m_0+1)$ 次连续可微。由上述条件以及第26.1.5节的讨论可见，带宽 $h_n$ 的收敛对这里2-TMLE的渐近行为起着关键作用。

需要指出，若 $K$ 确实取为区间 $(-0.5,+0.5)$ 上均匀核的张量积，则 $m_0=1$，条件2会变得更严格。当 $d\le2$ 时，确实存在使两个条件都成立的带宽速率 $h_n$；但当 $d>2$ 时并非如此。因此，尽管高阶核会增大变差范数，进而对 $h_n$ 施加更严格的条件，但当潜在混杂因素向量包含两个以上分量时，为控制表示误差，仍必须使用高阶核，以充分利用 $(A,W)$ 分布的底层光滑性。

与相应1-TMLE相比，2-TMLE通常可以在更大的模型中渐近线性且有效。一方面，通常只要1-TMLE有效，2-TMLE也就有效。这意味着2-TMLE处于一个安全范围内：我们预期，为构造2-TMLE而增加的靶向处理，不会损害一个渐近有效的1-TMLE。另一方面，在1-TMLE并不有效的许多情形下，2-TMLE仍然有效。例如，假设 $W$ 是具有足够光滑密度函数的一元随机变量。再假设 $\bar g_0$ 足够光滑，使一元二阶核平滑器能够产生 $\bar g_0$ 的最优估计量。此时，1-TMLE的有效性要求 $\bar Q_n$ 以快于 $n^{-1/10}$ 的速率趋于 $Q_0$。相比之下，相应2-TMLE所依据的二阶典范梯度，若用最优二阶核平滑器近似，则只要 $\bar Q_n$ 是 $\bar Q_0$ 的一致估计量，就能保证有效性，无论实际收敛速率是多少。如果 $W$ 有多个分量，这一问题会更加严重。例如，当 $W$ 为五维时，1-TMLE要求 $\bar Q_n$ 以快于 $n^{-5/18}$ 的速率趋于 $\bar Q_0$；相应2-TMLE基于二阶典范梯度的三阶核平滑近似，只要求 $\bar Q_n$ 以快于 $n^{-1/5}$ 的速率趋于 $\bar Q_0$。使用最优二阶核平滑器可以实现后一个速率，却无法实现前一个速率；如果不对 $\bar Q_0$ 施加进一步的光滑性假设，1-TMLE通常不会有效。

我们可能希望，通过把 $D_h^{(2)}$ 定义中的 $K_h(w_1-w_2)$，替换为基于 $w_1$ 与 $w_2$ 处倾向评分的核差异度，来系统地降维；Díaz 等（2016）讨论了这一点。如果 $\bar g_0$ 已知，就能够这样做。此时，可以使用 $K_h(\bar g_0(w_1)-\bar g_0(w_2))$ 定义降维后的近似二阶部分典范梯度，而不损失相应展开中余项的阶数。然而，在这一核差异度中，以 $\bar g$ 替代 $\bar g_0$，会在余项中引入一个二阶项，因而使使用这一二阶部分典范梯度的理论依据失效。不过，这并不排除沿这一路径获得有限样本收益的可能。事实上，下一小节将给出模拟研究的数值结果，在一个特定数据生成机制下，展示这一估计量的改进；我们将其记为1$^\star$-TMLE。

\subsection{数值结果}

我们给出一项模拟研究，展示1$^\star$-TMLE和2-TMLE相较于1-TMLE的表现，协变量维数为 $d=3$，样本量为 $n\in\{500,1000,2000,10{,}000\}$。这一模拟研究最初发表于 Díaz 等（2016）。对每个样本量 $n\in\{500,1000,2000,10{,}000\}$，我们从下列条件分布所蕴含的联合分布中，模拟1000个数据集：
\[
\begin{aligned}
W_1&\sim\operatorname{Beta}(2,2),\\
W_2\mid W_1=w_1&\sim\operatorname{Beta}(2w_1,2),\\
W_3\mid W_1=w_1,W_2=w_2&\sim\operatorname{Beta}(2w_1,2w_2),\\
A\mid W=(w_1,w_2,w_3)&\sim\operatorname{Bernoulli}\bigl(\expit(1+0.12w_1+0.1w_2+0.5w_3)\bigr),\\
Y\mid A=1,W=(w_1,w_2,w_3)&\sim\operatorname{Bernoulli}\bigl(\expit(-4+0.2w_1+0.3w_2+0.5\exp(w_3))\bigr).
\end{aligned}
\]
其中 $\operatorname{Bernoulli}(\pi)$ 表示成功概率为 $\pi$ 的Bernoulli分布，$\expit$ 表示函数 $u\mapsto\exp(u)/[1+\exp(u)]$，$\operatorname{Beta}(a,b)$ 表示参数为 $a$ 和 $b$ 的Beta分布。

对 $\bar Q_0$ 和 $g_0$ 的估计如下进行。对每个数据集，先拟合设定正确的参数模型。然后，对扰动参数 $p$，把 $\bar Q_n$ 的线性预测子乘以一个服从 $U(1-n^{-p},1)$ 的随机变量，再减去一个均值为 $3\times n^{-p}$、标准差为 $n^{-p}$ 的Gaussian随机变量。同样，对扰动参数 $q$，把 $g_n$ 的线性预测子乘以一个服从 $U(1-n^{-q},1)$ 的随机变量，再减去一个均值为 $3\times n^{-q}$、标准差为 $n^{-q}$ 的Gaussian随机变量。$p$ 和 $q$ 均在 $\{0.01,0.02,0.05,0.1,0.2,0.5\}$ 中变化。

在设定正确的参数模型中对MLE作上述扰动，是为了得到具有不同收敛速率的初始估计量。例如，假设 $g_{n,\mathrm{MLE}}$ 表示正确参数模型中 $g_0$ 的MLE，以 $g_{n,\mathrm{MLE},q}$ 表示扰动后的估计量。令 $U_n$ 和 $V_n$ 分别服从 $U(1-n^{-q},1)$ 和 $N(-3n^{-q},n^{-2q})$。将 $g_{n,\mathrm{MLE},q}:=U_ng_{n,\mathrm{MLE}}+V_n$ 代入 $\|\hat g_q^{\mathrm{MLE}}-g_0\|_{P_0}^2$，得到
\[
\begin{aligned}
\|g_{n,\mathrm{MLE},q}-g_0\|_{P_0}^2
&\le\|U_n(g_{n,\mathrm{MLE},q}-g_0)\|_{P_0}^2\\
&\quad+\|g_0(U_n-1)\|_{P_0}^2+\|V_n\|_{P_0}^2\\
&=O_P(n^{-1}+n^{-2q}).
\end{aligned}
\]
现在考虑不同的 $q$ 值。例如，$q=0.5$ 给出参数一致性速率 $\|g_{n,\mathrm{MLE},q}-g_0\|_{P_0}^2=O_P(n^{-1})$，而 $q=0$ 则导致不一致。

对通过这一扰动得到的每个初始估计量 $(\hat{\bar Q},\hat g)$，我们计算1-TMLE、1$^\star$-TMLE和2-TMLE。通过放大 $n^{1/2}$ 倍的偏差、相对于非参数效率界的相对方差，以及基于真实方差的95\%置信区间覆盖概率，比较这些估计量的表现。使用真实方差，可以将估计量本身的表现与其方差估计量的表现分开考察。真实抽样方差、偏差和覆盖概率，均通过1000次模拟运行的经验均值近似。

表26.1对扰动参数 $(p,q)$ 的部分取值，给出相对方差（rVar，定义为方差乘以 $n$ 后除以效率界）、放大 $n^{1/2}$ 倍的绝对偏差，以及使用估计量真实方差时95\%置信区间的覆盖概率。图26.1给出各估计量乘以 $n^{1/2}$ 的绝对偏差，图26.2给出模拟中使用的所有 $(p,q)$ 取值下95\%置信区间的覆盖概率。

我们注意到，对于某些缓慢的收敛速率，例如 $(p,q)=(0.01,0.01)$ 或 $(p,q)=(0.1,0.01)$，所有估计量的偏差都很大。相比之下，对于另一些缓慢的收敛速率，例如 $(p,q)=(0.01,0.1)$，1-TMLE经 $n^{1/2}$ 缩放后的绝对偏差，比2-TMLE和1$^\star$-TMLE发散得快得多。所提出估计量在大样本绝对偏差上的改进，在某些小样本情形中，是以方差增大为代价的。例如，在 $n\le2000$ 时，若结局模型收敛得足够快，而缺失机制并非如此，如 $(p,q)=(0.5,0.1)$，1-TMLE的方差低于其他估计量。正如理论所预测的，一阶TMLE的这一优势会渐近消失。

就覆盖而言，当 $p$ 和 $q$ 都较小时，使用1$^\star$-TMLE和2-TMLE带来的改进十分明显。例如，当 $n=2000$ 且 $(p,q)=(0.01,0.1)$ 时，1-TMLE、1$^\star$-TMLE和2-TMLE的估计覆盖概率分别为0、0.9和0.8。这一模拟展示了使用1$^\star$-TMLE和2-TMLE可能获得的巨大改进；但其代价是，当收敛速率足够快时，例如 $n\le2000$ 且 $(p,q)=(0.5,0.1)$，在小样本中会出现过度覆盖。

% 表26.1：27×12=324数据格；p栏共享三行，保留原空白。
\begingroup
\footnotesize
\setlength{\tabcolsep}{3pt}
\renewcommand{\arraystretch}{1.22}
\begin{longtable}{@{}ccrrrr|rrrr|rrrr@{}}
\caption{估计量表现随样本量及$\bar Q_0$、$g_0$初始估计量收敛速率的变化}\label{tab:ch26-performance}\\
\toprule
& & \multicolumn{4}{c}{1-TMLE} & \multicolumn{4}{c}{1$^\star$-TMLE} & \multicolumn{4}{c}{2-TMLE}\\
\cmidrule(lr){3-6}\cmidrule(lr){7-10}\cmidrule(l){11-14}
& & \multicolumn{4}{c}{$n$} & \multicolumn{4}{c}{$n$} & \multicolumn{4}{c}{$n$}\\
$p$ & $q$ & 500 & 1000 & 2000 & $10{,}000$ & 500 & 1000 & 2000 & $10{,}000$ & 500 & 1000 & 2000 & $10{,}000$\\
\midrule
\endfirsthead
\multicolumn{14}{l}{\footnotesize 表26.1（续）}\\
\toprule
& & \multicolumn{4}{c}{1-TMLE} & \multicolumn{4}{c}{1$^\star$-TMLE} & \multicolumn{4}{c}{2-TMLE}\\
\cmidrule(lr){3-6}\cmidrule(lr){7-10}\cmidrule(l){11-14}
& & \multicolumn{4}{c}{$n$} & \multicolumn{4}{c}{$n$} & \multicolumn{4}{c}{$n$}\\
$p$ & $q$ & 500 & 1000 & 2000 & $10{,}000$ & 500 & 1000 & 2000 & $10{,}000$ & 500 & 1000 & 2000 & $10{,}000$\\
\midrule
\endhead
\midrule
\multicolumn{14}{r}{\footnotesize 续下页}\\
\endfoot
\bottomrule
\endlastfoot
\multicolumn{14}{l}{\textbf{$n^{1/2}|\text{偏差}|$}}\\*
 & 0.01 & 3.02 & 4.34 & 6.14 & 13.56 & 1.39 & 1.97 & 2.77 & 5.98 & 0.47 & 1.00 & 2.69 & 6.40\\*
0.01 & 0.10 & 1.93 & 2.55 & 3.27 & 5.65 & 0.18 & 0.22 & 0.32 & 0.41 & 1.47 & 1.95 & 0.61 & 0.73\\*
 & 0.50 & 0.07 & 0.05 & 0.04 & 0.03 & 0.09 & 0.09 & 0.12 & 0.15 & 1.71 & 2.05 & 0.67 & 0.63\\*
\cmidrule(l){1-14}
 & 0.01 & 1.33 & 1.77 & 2.31 & 4.26 & 0.63 & 0.84 & 1.17 & 2.22 & 0.03 & 0.28 & 1.08 & 2.28\\*
0.10 & 0.10 & 0.87 & 1.05 & 1.25 & 1.70 & 0.08 & 0.11 & 0.14 & 0.17 & 0.96 & 1.02 & 0.22 & 0.12\\*
 & 0.50 & 0.01 & 0.02 & 0.01 & 0.03 & 0.07 & 0.05 & 0.07 & 0.04 & 0.87 & 0.90 & 0.21 & 0.23\\*
\cmidrule(l){1-14}
 & 0.01 & 0.09 & 0.08 & 0.08 & 0.07 & 0.00 & 0.01 & 0.03 & 0.03 & 0.02 & 0.02 & 0.04 & 0.02\\*
0.50 & 0.10 & 0.04 & 0.03 & 0.03 & 0.00 & 0.00 & 0.07 & 0.19 & 0.01 & 0.02 & 0.09 & 0.11 & 0.03\\*
 & 0.50 & 0.01 & 0.00 & 0.01 & 0.01 & 0.02 & 0.00 & 0.01 & 0.01 & 0.15 & 0.08 & 0.01 & 0.03\\
\midrule
\multicolumn{14}{l}{\textbf{相对方差（rVar）}}\\*
 & 0.01 & 1.60 & 1.59 & 1.57 & 1.60 & 3.22 & 2.13 & 2.15 & 2.11 & 2.58 & 2.87 & 1.89 & 2.09\\*
0.01 & 0.10 & 1.73 & 1.72 & 1.56 & 1.46 & 1.05 & 1.09 & 1.02 & 0.99 & 2.14 & 1.97 & 1.27 & 1.17\\*
 & 0.50 & 1.06 & 1.08 & 1.03 & 1.01 & 1.10 & 1.08 & 0.97 & 1.07 & 2.15 & 2.10 & 1.31 & 1.17\\*
\cmidrule(l){1-14}
 & 0.01 & 1.10 & 1.07 & 1.06 & 1.04 & 1.13 & 1.26 & 1.14 & 1.25 & 1.71 & 1.58 & 1.31 & 1.19\\*
0.10 & 0.10 & 1.21 & 1.20 & 1.09 & 1.04 & 0.95 & 0.99 & 0.99 & 0.92 & 1.75 & 1.78 & 1.17 & 1.06\\*
 & 0.50 & 0.97 & 0.98 & 1.02 & 0.98 & 1.03 & 1.01 & 1.01 & 1.01 & 2.17 & 1.96 & 1.19 & 1.04\\*
\cmidrule(l){1-14}
 & 0.01 & 1.02 & 1.04 & 1.01 & 1.01 & 1.04 & 1.00 & 1.05 & 1.03 & 1.02 & 1.02 & 0.99 & 0.99\\*
0.50 & 0.10 & 1.04 & 0.95 & 0.97 & 0.97 & 1.14 & 6.10 & 17.17 & 1.02 & 3.58 & 9.19 & 16.22 & 0.84\\*
 & 0.50 & 0.99 & 0.98 & 1.01 & 0.95 & 1.10 & 0.93 & 1.00 & 0.93 & 1.37 & 1.23 & 1.10 & 0.97\\
\midrule
\multicolumn{14}{l}{\textbf{覆盖概率（Cov. P）}}\\*
 & 0.01 & 0.01 & 0.00 & 0.00 & 0.00 & 0.76 & 0.23 & 0.03 & 0.00 & 0.91 & 0.78 & 0.03 & 0.00\\*
0.01 & 0.10 & 0.19 & 0.03 & 0.00 & 0.00 & 0.93 & 0.92 & 0.90 & 0.86 & 0.49 & 0.22 & 0.80 & 0.73\\*
 & 0.50 & 0.94 & 0.95 & 0.94 & 0.94 & 0.93 & 0.95 & 0.95 & 0.94 & 0.38 & 0.22 & 0.80 & 0.79\\*
\cmidrule(l){1-14}
 & 0.01 & 0.30 & 0.08 & 0.01 & 0.00 & 0.79 & 0.69 & 0.42 & 0.03 & 0.95 & 0.93 & 0.54 & 0.02\\*
0.10 & 0.10 & 0.66 & 0.53 & 0.35 & 0.10 & 0.95 & 0.94 & 0.94 & 0.93 & 0.70 & 0.67 & 0.93 & 0.94\\*
 & 0.50 & 0.95 & 0.95 & 0.95 & 0.94 & 0.94 & 0.95 & 0.94 & 0.94 & 0.79 & 0.77 & 0.92 & 0.93\\*
\cmidrule(l){1-14}
 & 0.01 & 0.95 & 0.95 & 0.95 & 0.95 & 0.95 & 0.95 & 0.96 & 0.96 & 0.95 & 0.95 & 0.94 & 0.96\\*
0.50 & 0.10 & 0.94 & 0.94 & 0.95 & 0.95 & 0.95 & 0.99 & 0.98 & 0.95 & 0.99 & 0.98 & 0.98 & 0.95\\*
 & 0.50 & 0.95 & 0.95 & 0.95 & 0.95 & 0.94 & 0.95 & 0.95 & 0.96 & 0.93 & 0.95 & 0.95 & 0.95\\
\end{longtable}
\endgroup


在小样本中，只要其中一个模型收敛得足够快，即 $p$ 或 $q$ 中任一个等于0.5，2-TMLE的偏差表现就比1-TMLE和1$^\star$-TMLE更差，进而导致较低的覆盖概率。随着 $n$ 增大，这一问题有所减轻。虽然这提示在小样本中使用2-TMLE可能需要谨慎，但应注意，实践中使 $p$ 或 $q$ 达到0.5，对应于使用设定正确的参数模型；在大多数应用中，这几乎不可能做到。

图26.1和图26.2显示，存在一个收敛速率较慢的区域，所提出的估计量在该区域中优于标准一阶TMLE。此外，如图26.1所示，在偏差方面，2-TMLE相对于1$^\star$-TMLE有一点优势。不过，总体而言，我们没有看到2-TMLE相对于1$^\star$-TMLE的实际优势。事实上，在小样本中1$^\star$-TMLE优于2-TMLE，在所有样本量下都优于1-TMLE，但仍需注意前述某些情形下方差增大的问题。

\begin{figure}[p]
\centering
\resizebox{.91\textwidth}{!}{% 原书真实矢量/位图数据与全部坐标刻度保留，288dpi整块逐像素相同。
% 仅在十二面板与色条之外覆盖英文图例标题并叠加中文。
\begin{tikzpicture}[x=1bp,y=-1bp]
\node[anchor=north west,inner sep=0pt] at (0,0) {\includegraphics[width=269bp,height=258bp]{figures/ch26_fig1_data.pdf}};
\path[fill=white,draw=none] (240,100.5) rectangle (268,110.6);
\node[anchor=west,inner sep=0pt,font=\fontsize{7bp}{8bp}\selectfont] at (240.6,106.4) {$\sqrt n$\,偏差};
\end{tikzpicture}
}
\caption{经$n^{1/2}$缩放的绝对偏差随样本量及$\bar Q_0$、$g_0$初始估计量收敛速率的变化。}
\label{fig:ch26-bias}
\end{figure}

\begin{figure}[p]
\centering
\resizebox{.91\textwidth}{!}{% 原书真实矢量/位图数据与全部坐标刻度保留，288dpi整块逐像素相同。
% 仅在十二面板与色条之外覆盖英文图例标题并叠加中文。
\begin{tikzpicture}[x=1bp,y=-1bp]
\node[anchor=north west,inner sep=0pt] at (0,0) {\includegraphics[width=269bp,height=258bp]{figures/ch26_fig2_data.pdf}};
\path[fill=white,draw=none] (240.5,99) rectangle (267,107.6);
\node[anchor=west,inner sep=0pt,font=\fontsize{6.4bp}{7.2bp}\selectfont] at (241,103.3) {覆盖概率};
\end{tikzpicture}
}
\caption{置信区间覆盖概率随样本量及$\bar Q_0$、$g_0$初始估计量收敛速率的变化。}
\label{fig:ch26-coverage}
\end{figure}
\FloatBarrier
